\section{Operator Fredholm}
\begin{defn} Misalkan $H$ adalah ruang Hilbert dan $\ofml K(H)$ merupakan ideal operator kompak di dalam
aljabar-$C^*$~$\ofml B(H)$. Maka aljabar hasil bagi (lihat proposisi~\ref{001902})
$\ofml Q(H) := \ofml B(H)/\ofml K(H)$ adalah
 \index{aljabar Calkin}%
 \index{aljabar!Calkin}%
 \index{Q@$\ofml Q(H)$ (aljabar Calkin)}%
\df{aljabar Calkin}. Seperti biasa, pemetaan hasil bagi dari $\ofml B(H)$ ke $\ofml
B(H)/\ofml K(H)$
 \index{pi@$\pi$ (pemetaan hasil bagi)}%
 \index{pi@$\pi(T)$ (elemen aljabar Calkin yang memuat $T$)}%
dilambangkan dengan $\pi$, sehingga jika $T \in \ofml B(H)$ maka $\pi(T) = [T] = T + \ofml K(H)$ adalah
elemen yang bersesuaian dengannya dalam aljabar Calkin. Suatu elemen $T \in \ofml B(H)$ adalah
 \index{Fredholm!operator}%
 \index{operator!Fredholm}%
\df{operator Fredholm} jika $\pi(T)$ invertibel di dalam~$\ofml Q(H)$. Keluarga
 \index{F@$\ofml F(H)$ (operator Fredholm pada $H$)}%
semua operator Fredholm pada $H$ dilambangkan dengan~$\ofml F(H)$.
\end{defn}

\begin{prop} Jika $H$ ruang Hilbert, maka $\ofml F(H)$ merupakan subhimpunan terbuka swaadjoin
dari~$\ofml B(H)$ dan tertutup terhadap
 \index{perturbasi}%
perturbasi kompak (yakni, jika $T$ Fredholm dan $K$ kompak, maka $T + K$ juga Fredholm).
\end{prop}

\begin{thm}[Teorema Atkinson]\label{004352}
 \index{teorema Atkinson}%
Suatu operator ruang Hilbert bersifat Fredholm jika dan hanya jika kernel dan kokernelnya
berdimensi hingga.
\end{thm}

\begin{proof} Lihat \cite{Arveson:2002}, teorema 3.3.2; \cite{Blackadar:2006}, teorema I.8.3.6;
\cite{Douglas:1972}, teorema 5.17; \cite{HigsonR:2000}, teorema 2.1.4;
atau~\cite{Wegge-Olsen:1993}, teorema 14.1.1. \ns
\end{proof}

Menariknya, dimensi kernel dan kokernel suatu operator Fredholm bukanlah besaran yang
sangat penting. Namun, selisih keduanya ternyata sangat penting.

\begin{defn} Jika $T$ adalah operator Fredholm pada suatu ruang Hilbert, maka
 \index{indeks!operator Fredholm}%
 \index{Fredholm!indeks (\seeonly{indeks})}%
\df{indeks Fredholm} operator tersebut (atau cukup \df{indeks}) didefinisikan oleh
   \[ \ind T := \dim\ker T - \dim\coker T\,. \]
\end{defn}

\begin{exam}\label{0044125} Setiap operator ruang Hilbert yang invertibel bersifat Fredholm
 \index{indeks!operator invertibel}%
 \index{operator!invertibel!indeks}%
 \index{invertibel!operator!indeks}%
dengan indeks nol.
\end{exam}

\begin{exam}\label{0044128} Indeks Fredholm dari
 \index{indeks!operator geser unilateral}%
 \index{operator geser unilateral!indeks}%
 \index{operator!geser unilateral!indeks}%
operator geser unilateral adalah~$-1$.
\end{exam}

\begin{exam} Untuk pemetaan antar-ruang, kita memakai konvensi standar: pemetaan linear disebut Fredholm jika
jangkauannya tertutup serta kernel dan kokernelnya berdimensi hingga. Dengan konvensi ini, setiap pemetaan linear $T\colon V \sto W$ antara ruang-ruang vektor berdimensi hingga bersifat Fredholm dan
$\ind T = \dim V - \dim W$. \end{exam}

\begin{exam}\label{004413} Jika $T$ adalah operator Fredholm pada suatu ruang Hilbert,
 \index{indeks!adjoin operator}%
 \index{operator!adjoin!indeks}%
 \index{adjoin!indeks}%
maka $\ind T^* = -\ind T$.
\end{exam}

\begin{exam} Indeks setiap operator Fredholm
 \index{indeks!operator Fredholm normal}%
 \index{operator!Fredholm normal!indeks}%
 \index{normal!operator Fredholm!indeks}%
normal adalah~$0$.
\end{exam}

\begin{lem}\label{004432} Misalkan $S\colon U \sto V$ dan $T\colon V \sto W$ adalah transformasi linear
antara ruang-ruang vektor. Maka (terdapat pemetaan linear sedemikian sehingga) barisan
berikut eksak.
 \[ \vc 0 \to \ker S \to \ker TS \to \ker T \to \coker S
                            \to \coker TS \to \coker T \to \vc 0\,. \]
\end{lem}

\begin{lem}\label{004433} Jika $V_0, V_1, \dots, V_n$ adalah ruang-ruang vektor berdimensi hingga
dan barisan
  \[ \vc 0 \to V_0 \to V_1 \to \dots \to V_n \to \vc 0 \]
eksak, maka
  \[ \sum_{k=0}^n (-1)^k \dim V_k = 0\,. \]
\end{lem}

\begin{prop} Misalkan $H$ ruang Hilbert berdimensi tak hingga. Maka himpunan $\ofml F(H)$ yang terdiri atas operator Fredholm
pada $H$ merupakan semigrup terhadap komposisi dan fungsi indeks $\ind$ adalah epimorfisme
dari $\ofml F(H)$ ke semigrup aditif bilangan bulat~$\Z$.
\end{prop}

\begin{proof} \emph{Petunjuk.} Gunakan \ref{004432}, \ref{004433}, \ref{0044125}, \ref{0044128},
dan~\ref{004413}. \ns
\end{proof}





















\section{Alternatif Fredholm -- Penutup}
Hasil berikut menyiratkan bahwa setiap operator Fredholm berindeks nol dapat ditulis sebagai
jumlah suatu operator invertibel dan operator kompak.

\begin{prop}\label{004712} Jika $T$ adalah operator Fredholm berindeks nol pada suatu ruang Hilbert,
maka terdapat isometri parsial berperingkat hingga $V$ sedemikian sehingga $T - V$ bersifat invertibel.
\end{prop}

\begin{proof} Lihat \cite{Pedersen:1995}, lemma 3.3.14 atau \cite{Wegge-Olsen:1993}, proposisi
14.1.3.  \ns
\end{proof}

\begin{lem} Jika $F$ adalah operator berperingkat hingga pada suatu ruang Hilbert, maka $I + F$ bersifat Fredholm dengan
indeks nol.
\end{lem}

\begin{proof} Lihat \cite{Pedersen:1995}, lemma 3.3.13 atau \cite{Wegge-Olsen:1993}, lemma 14.1.4. \ns
\end{proof}

\begin{notn} Misalkan $H$ adalah ruang Hilbert. Untuk setiap bilangan bulat $n$, kita melambangkan
 \index{F@$\ofml F_n(H)$, $\ofml F_n$ (operator Fredholm berindeks~$n$)}%
keluarga semua operator Fredholm berindeks $n$ pada $H$ dengan $\ofml F_n(H)$ atau cukup $\ofml
F_n$.
\end{notn}

\begin{prop}\label{004724} Untuk setiap ruang Hilbert berlaku $\ofml F_0 + \ofml K = \ofml F_0$.
\end{prop}

\begin{proof} Lihat \cite{Douglas:1972}, lemma 5.20 atau \cite{Wegge-Olsen:1993}, 14.1.5. \ns
\end{proof}

\begin{cor}[Alternatif Fredholm V]\label{004725}
 \index{alternatif Fredholm!versi V}%
Setiap operator Riesz--Schauder pada suatu ruang Hilbert bersifat Fredholm berindeks nol.
\end{cor}

\begin{proof} Ini segera mengikuti proposisi sebelumnya~\ref{004724} dan
contoh~\ref{0044125}.
\end{proof}

Kita sebenarnya telah membuktikan hasil yang lebih kuat: versi terakhir (yang cukup umum) dari
\emph{alternatif Fredholm}.

\begin{cor}[Alternatif Fredholm VI]\label{0047253} Suatu operator ruang Hilbert merupakan operator
 \index{alternatif Fredholm!versi VI}%
Riesz--Schauder jika dan hanya jika operator itu merupakan operator Fredholm berindeks nol.
\end{cor}

\begin{proof} Proposisi~\ref{004712} dan korolari~\ref{004725}.
\end{proof}

\begin{prop} Jika $T$ adalah operator Fredholm berindeks $n \in \Z$ pada ruang Hilbert $H$
dan $K$ operator kompak pada~$H$, maka $T + K$ juga merupakan operator Fredholm berindeks~$n$; yakni,
  \[ \ofml F_n + \ofml K = \ofml F_n\,. \]
\end{prop}

\begin{proof} Lihat \cite{HigsonR:2000}, proposisi 2.1.6; \cite{Pedersen:1995}, teorema
3.3.17; atau \cite{Wegge-Olsen:1993}, proposisi 14.1.6. \ns
\end{proof}

\begin{prop} Misalkan $H$ adalah ruang Hilbert. Maka $\ofml F_n(H)$ merupakan subhimpunan terbuka dari $\ofml B(H)$
untuk setiap bilangan bulat~$n$.
\end{prop}

\begin{proof} Lihat \cite{Wegge-Olsen:1993}, proposisi 14.1.8. \ns \end{proof}

\begin{defn}\label{004735} Suatu
 \index{lintasan}%
\df{lintasan} dalam ruang topologis $X$ adalah pemetaan kontinu dari interval $[0,1]$
ke~$X$. Dua titik $p$ dan $q$ dalam $X$ dikatakan
 \index{lintasan!titik-titik yang terhubung oleh}%
 \index{terhubung!oleh lintasan}%
\df{terhubung oleh lintasan} (atau
 \index{homotop!titik-titik dalam ruang topologis}%
\df{homotop dalam~$X$}) jika terdapat lintasan $f\colon [0,1] \sto X$ dalam $X$ sedemikian sehingga
$f(0) = p$ dan $f(1) = q$. Dalam hal ini kita
 \index{<binrelequivh@$p \sim_h q$ (homotopi titik-titik)}%
menulis $p \sim_h q$ dalam~$X$ (atau cukup $p \sim_h q$ ketika ruang $X$ jelas dari konteks).
\end{defn}


\begin{exam} Misalkan $a$ adalah elemen invertibel dalam aljabar-$C^*$ beridentitas $A$ dan $b$ adalah elemen $A$
sedemikian sehingga $\norm{a - b} < \norm{a^{-1}}^{-1}$. Maka $a \sim_h b$ dalam~$\inv(A)$.
\end{exam}

\begin{proof}[\emph{Petunjuk bukti}] Terapkan korolari~\ref{cor2_Neumann_series} pada titik-titik dalam ruas garis tertutup~$[a,b]$. \ns
\end{proof}

\begin{prop}\label{K002011} Relasi $\sim_h$ untuk ekuivalensi homotopi yang didefinisikan di atas merupakan suatu relasi
 \index{ekuivalensi!homotopi!titik-titik}%
ekuivalensi pada himpunan titik-titik suatu ruang topologis.
\end{prop}

\begin{defn} Jika $X$ adalah ruang topologis dan $\sim_h$ adalah relasi ekuivalensi homotopi,
maka kelas-kelas ekuivalensi yang dihasilkan adalah
 \index{lintasan!komponen}%
 \index{komponen!lintasan}%
\df{komponen lintasan} dari~$X$.
\end{defn}

Proposisi berikut mengidentifikasi komponen-komponen lintasan himpunan operator Fredholm sebagai
himpunan $\ofml F_n$ yang terdiri atas operator berindeks~$n$.

\begin{prop} Operator $S$ dan $T$ di dalam ruang operator Fredholm $\ofml F(H)$ pada
ruang Hilbert $H$ bersifat homotop dalam $\ofml F(H)$ jika dan hanya jika keduanya memiliki indeks yang sama.
\end{prop}

\begin{proof} Lihat \cite{Wegge-Olsen:1993}, korolari 14.1.9. \ns \end{proof}







\endinput
